% ============================================================
% One-page progress summary -- compiles as-is on Overleaf.
% If the document exceeds one page, change 10pt to 9pt.
% ============================================================

\documentclass[10pt,conference]{IEEEtran}

\usepackage[T1]{fontenc}
\usepackage[utf8]{inputenc}
\usepackage{amsmath}
\usepackage{amssymb}
\usepackage{booktabs}
\usepackage{float}
\usepackage{enumitem}

\setlist{
    noitemsep,
    topsep=2pt,
    partopsep=0pt,
    leftmargin=*
}

\begin{document}

\title{Progress Summary}

\author{\IEEEauthorblockN{Jad Al Lakkis}}

\maketitle
\thispagestyle{empty}

\begin{abstract}
Progress since the last meeting: a working constrained-PPO baseline for tile-based 360-degree video streaming, implemented, trained, and validated, with measured results.
\end{abstract}

\section{Implementation}

\noindent\textbf{Environment.} Gymnasium environment for the Runner video: 12 valid head-movement traces, 64 tiles/frame (8$\times$8), 30 fps, 30-frame GOPs = 1 s/slot. Integrates the Rician-fading channel model, buffer/stall dynamics, viewport-coverage computation, and the per-tile rate-distortion (QP-level) layer model into one step function.

\smallskip
\noindent\textbf{State/action.} $\omega^t=(Z^t,\Delta_\theta^{t-1},\Delta_\phi^{t-1},h^t)$: buffer level, previous-step prediction error (both angles), channel gain. Action: 64-tile enhance mask + 1 of 5 discrete transmit-power levels (0/25/50/75/100\% of $P_{\max}$). Tiles overlapping the predicted viewport are enhanced unconditionally (eq.~2); the agent only chooses among the rest.

\smallskip
\noindent\textbf{Constraint calibration.} $\bar D,\bar P,\bar B$ derived experimentally, not assumed: 3 reference tile policies (mandatory-only / +50\% random / all tiles) $\times$ 3 power levels, run on the 9 training traces with fixed seeds, measuring the exact discounted $J_D,J_P,J_B$. Values chosen inside the resulting efficient/wasteful range (Table~\ref{tab:budgets}).

\smallskip
\noindent\textbf{Training.} PPO (Stable-Baselines3), 2{,}048-step rollouts ($\approx$57 episodes each, each episode 36 steps). Every step uses the Lagrangian reward
\[
r_{\mathrm L}^t=\beta_q\bar Q^t_{\text{viewport}}-\mu_D\frac{D^t}{T_0}-\mu_P\frac{P^t}{P_{\max}}-\mu_B\frac{\sum_i x_i^t}{N},
\]
with $\mu_D,\mu_P,\mu_B$ updated once per rollout by projected dual ascent on that rollout's average constraint costs:
\[
\mu_j \leftarrow \max\!\Bigl(0,\ \mu_j+\eta\bigl(\overline{J_j}^{\,\text{rollout}}-\bar b_j\bigr)\Bigr),\ \ j\in\{D,P,B\}.
\]
Observations normalized online (VecNormalize), fit on training traces only, frozen at validation. 3 of 12 traces held out for validation with fixed seeds, evaluated every rollout -- results in Tables~\ref{tab:startend}--\ref{tab:rollouts}. Best checkpoint kept by feasibility first, then highest $J_Q$.

\section{Formulation}
\[
\max_\pi J_Q \quad \text{s.t.} \quad J_D\le\bar D,\ J_P\le\bar P,\ J_B\le\bar B
\]
Solved via Lagrangian relaxation; the resulting reward and multiplier update are given in the Training paragraph above. $\bar D,\bar P,\bar B$ enter only the update, not the per-step reward directly.

\begin{table}[H]
\centering
\footnotesize
\begin{tabular}{lccc}
\toprule
Budget & Efficient & Wasteful & Chosen \\
\midrule
$\bar D$ (stall) & $\approx0.00$ & $\approx0.60$ & $\mathbf{0.10}$ \\
$\bar P$ (power) & $0.00$--$0.505$ (min/mid) & $1.01$ (max) & $\mathbf{0.60}$ \\
$\bar B$ (tiles) & $\approx0.23$ & $\approx1.01$ & $\mathbf{0.35}$ \\
\bottomrule
\end{tabular}
\caption{Calibrated constraint budgets.}
\label{tab:budgets}
\end{table}

\section{Training and Results}
Runs of 50K, 300K, and 700K steps. Validation feasibility: $0/25\to0/147\to140/342$, with the final 20/20 evaluations feasible in the 700K run. The best-feasible checkpoint is not always the last one: a 300K-run checkpoint at 294{,}912 steps had lower violation ($\approx0.057$) than that run's own final step.

\begin{table}[H]
\centering
\footnotesize
\begin{tabular}{lccc}
\toprule
\textbf{Metric (700K run)} & \textbf{Start} & \textbf{End} & \textbf{$\Delta$} \\
\midrule
Feasible & No & Yes & -- \\
Coverage & 0.927 & 0.882 & $-4.9\%$ \\
Stall (s) & 0.531 & 0.022 & $-95.9\%$ \\
Power (mW) & 49.6 & 27.6 & $-44.4\%$ \\
Enhanced tiles & 39.1 & 19.0 & $-51.4\%$ \\
\bottomrule
\end{tabular}
\caption{Physical training metrics, start vs.\ end of the 700K run.}
\label{tab:startend}
\end{table}

\noindent Multipliers: $\mu_D$ 1.00$\to$1.71, $\mu_B$ 1.00$\to$1.46 (still exerting pressure); $\mu_P$ 1.00$\to$0.009 (power constraint slack at $\bar P=0.6$). PPO's explained variance rose from $\approx0$ to $0.41$.

\begin{table}[H]
\centering
\scriptsize
\begin{tabular}{lccccc}
\toprule
Rollout & Feas. & Cov. & Stall(s) & Pwr(mW) & Tiles \\
\midrule
1   & 0 & 0.918 & 0.422 & 72.9 & 39.6 \\
2   & 0 & 0.926 & 0.602 & 33.3 & 36.6 \\
341 & 1 & 0.841 & 0.001 & 25.0 & 17.1 \\
342 & 1 & 0.840 & 0.001 & 25.0 & 17.0 \\
\bottomrule
\end{tabular}
\caption{First/last 2 of 342 validation rollouts, 700K run (step = rollout $\times$ 2{,}048). Full log: \texttt{runs/train\_700k/train\_700k\_key\_metrics.csv}.}
\label{tab:rollouts}
\end{table}

\section{Questions}
\begin{itemize}
    \item Keep $\lambda=0.01$, or change it -- recalibrate the budgets if so?
    \item Are these the right validation metrics to be tracking?
    \item Tighten $\bar P$, or narrow the power action range (e.g.\ 0--50\%, still 5 slots), given $\mu_P\to0$?
\end{itemize}

\section{Next Steps}
\begin{enumerate}
    \item Decide on further training length beyond 700K.
    \item Resolve the questions above.
    \item Compare PPO vs.\ PPO + PDS + VE under a matched protocol.
\end{enumerate}

\end{document}
